% ============================================================================
% บทที่ 1 — เริ่มต้นใช้งาน
% ============================================================================
\mychapter{เริ่มต้นใช้งาน}

\noindent เปิด Exam Scheduler จากที่อยู่ (URL) ที่ได้รับจากผู้ดูแลระบบ
ระบบเริ่มต้นเป็นภาษาไทย ใช้ปุ่ม \texttt{English} / \texttt{ไทย} ที่มุมขวาบน
เพื่อสลับภาษา

\section{เมนูหลักและสถานะของแต่ละหน้า}

เมนูหลักอยู่ทางด้านซ้าย ประกอบด้วยรายการต่อไปนี้ เมนูที่ยังใช้ไม่ได้จะเป็นสีจาง
และกดไม่ได้ จนกว่างานจะดำเนินถึงขั้นที่เปิดใช้หน้านั้น บนหน้าจอขนาดเล็ก
เมนูหลักจะย้ายไปอยู่ใต้ปุ่ม \texttt{อื่น ๆ} ที่มุมขวาบน

\begin{longtable}{L{4cm}L{9.5cm}}
\toprule
\textbf{เมนู} & \textbf{หน้าที่และเงื่อนไขการเปิด} \\
\midrule
ข้อมูลและการตั้งค่า & หน้าเริ่มต้นสำหรับนำเข้าไฟล์และตั้งค่าปฏิทิน \\
ตรวจสอบข้อมูล & ตรวจสอบข้อมูลก่อนจัดตาราง (เปิดได้หลังนำเข้า) \\
ภาพรวม & สรุปจำนวนวิชาที่จัดแล้วและยังไม่ได้จัด (เปิดได้หลังจัดตาราง) \\
รายวิชา & ค้นหาและแก้ไขเวลาของแต่ละวิชา (เปิดได้หลังจัดตาราง) \\
ปฏิทิน & แสดงตารางสอบรายสัปดาห์ (เปิดได้หลังจัดตาราง) \\
กลุ่มนักศึกษา & ดูตารางสอบของแต่ละกลุ่มนักศึกษา (เปิดได้หลังจัดตาราง) \\
ข้อควรตรวจสอบ & รายการปัญหาที่ตรวจพบ (เปิดได้หลังนำเข้า) \\
บันทึกและแชร์ & บันทึกขึ้นคลาวด์และสร้างลิงก์แชร์ \\
ตรวจสอบ & หน้ารายงานการตรวจสอบคุณภาพตาราง \\
วิธีใช้งาน คู่มือและตัวอย่าง & เปิดคู่มือ รูปแบบข้อมูล ตัวอย่าง และการแก้ปัญหา \\
ดาวน์โหลด project.zip & ดาวน์โหลดไฟล์เก็บถาวรครบชุด \\
\bottomrule
\end{longtable}

ชื่อของสามหน้าใกล้เคียงกันแต่หน้าที่ต่างกัน: \textbf{ตรวจสอบข้อมูล}
ใช้ตรวจข้อมูลที่นำเข้า \textbf{ข้อควรตรวจสอบ} แสดงรายการปัญหาที่ตรวจพบ
และ \textbf{ตรวจสอบ} ใช้ดูรายงานคุณภาพของตารางหลังจัด
(รายละเอียดอยู่ในบทที่ 6)

\section{เริ่มต้นอย่างรวดเร็วด้วยข้อมูลตัวอย่าง}

ระบบมีข้อมูลตัวอย่างในตัวเพื่อให้เห็นผลลัพธ์ทันทีโดยไม่ต้องเตรียมไฟล์เอง:

\begin{steps}
  \item ที่หน้าข้อมูลและการตั้งค่า ให้กดปุ่ม \texttt{ใช้ไฟล์ตัวอย่าง}
        (อยู่มุมขวาบนของส่วนตั้งค่า)
  \item กดปุ่ม \texttt{นำเข้าและตรวจสอบ} ด้านล่าง
  \item รอจนระบบนำเข้าเสร็จ แล้วกด \texttt{จัดตาราง}
\end{steps}

\begin{expected}
หน้าภาพรวมแสดงจำนวนวิชาที่จัดแล้วและยังไม่ได้จัด โดยข้อมูลตัวอย่างจะถูกจัดให้ครบ
\end{expected}

สามารถดูและคัดลอกตัวอย่างข้อมูลเพิ่มเติมได้จากหน้า
\textbf{วิธีใช้งาน คู่มือและตัวอย่าง}

\manualfigure{docs/usage-report/screenshots/figures/F02.png}{0.95\textwidth}{%
  หน้าหลักของระบบหลังกดใช้ไฟล์ตัวอย่าง}

\seealso{บทที่ 3 สำหรับการนำเข้าไฟล์จริง และบทที่ 2 สำหรับการเตรียมไฟล์ข้อมูล}
